% V.5: in a group project each participant describes their own contribution, at
% least two pages each. This is a two-author project, so this chapter is
% MANDATORY and must not be removed. It is also the counterpart of the explicit
% criterion in IV.6, "justification of the participation of the student", which
% sits inside the 60 % of the mark given to the work itself.
%
% Each student writes their own section, in the first person, and the two
% sections must be clearly different from one another. Each student also
% presents this material in English at the defence (VI.3).

% Which of the two authors is Student 1 and which is Student 2 is still to be
% agreed between them; the \todo below marks where the names go once that is
% settled. The two profiles
% themselves come from the supervisor's assignment: Student 1 works on
% representation learning and model evaluation, Student 2 on graph
% construction, fuzzy techniques and results analysis.

\chapter{Personal contributions}
\label{cap:contributions}

\section{\todo{Student 1 name} --- Representation learning and model evaluation}
\label{sec:contribution-first}

\todo{At least two pages. Write it in the first person and make it concrete.
Cover: the extractors (\texttt{extractors/*}: \texttt{mfcc}, \texttt{acoustic},
\texttt{gauss\_mfcc}, \texttt{posterior\_*}, \texttt{cnn},
\texttt{metric\_learning} for \texttt{siamese\_mel}/\texttt{triplet\_mel},
\texttt{vector\_siamese} for \texttt{siamese\_desc}/\texttt{siamese\_emb},
\texttt{mert}); the training code (\texttt{training/*}); the evaluation modules
owned by this line (\texttt{evaluation/\{retrieval, triplets, probes,
classifiers, representation\}}); ground truth (triplets, MTT facets, graded
genres); the classical classifiers of P9; experiments e04 and e05; and the
Methods and Evaluation sections on representations, siamese networks and ground
truth (also SQ1, SQ4 and SQ5 in the Results chapter). Which problems came up
and how they were solved. Point at specific files and configurations in the
repository (\repositorio) so that a claim can be checked.}

\section{\todo{Student 2 name} --- Graph construction, fuzzy techniques and results analysis}
\label{sec:contribution-second}

\todo{At least two pages, same instructions, and clearly distinct in content
from the section above. Cover: graph construction and comparison
(\texttt{graphs/*}), similarity measures (\texttt{similarity.py}), the fuzzy
block (\texttt{fuzzy/*}), the evaluation modules owned by this line
(\texttt{evaluation/\{structure, hubness, robustness, graph\_classification,
link\}}), plotting, experiments e03, e06, e07, e09, e10 and e11, and the
Methods sections on graphs, similarity measures and fuzzy modelling (also SQ2,
SQ3 and SQ6, and the Discussion chapter, in the Results/Discussion chapters).
Which problems came up and how they were solved. Point at specific files and
configurations in the repository (\repositorio) so that a claim can be
checked.}

\todo{Both authors: state jointly-owned work explicitly, since P1--P3 (the
baseline and the shared interface contract), P11 (model comparison, e08 and
e09, conclusions and the introduction) and this chapter itself are done
together.}
